%  LaTeX support: latex@mdpi.com 
%  For support, please attach all files needed for compiling as well as the log file, and specify your operating system, LaTeX version, and LaTeX editor.

%=================================================================
\documentclass[signals,article,submit,moreauthors,pdftex]{Definitions/mdpi} 

%--------------------
% Class Options:
%--------------------

%---------
% article
%---------
% The default type of manuscript is "article", but can be replaced by: 
% abstract, addendum, article, book, bookreview, briefreport, casereport, comment, commentary, communication, conferenceproceedings, correction, conferencereport, entry, expressionofconcern, extendedabstract, datadescriptor, editorial, essay, erratum, hypothesis, interestingimage, obituary, opinion, projectreport, reply, retraction, review, perspective, protocol, shortnote, studyprotocol, systematicreview, supfile, technicalnote, viewpoint, guidelines, registeredreport, tutorial
% supfile = supplementary materials

%----------
% submit
%----------
% The class option "submit" will be changed to "accept" by the Editorial Office when the paper is accepted. This will only make changes to the frontpage (e.g., the logo of the journal will get visible), the headings, and the copyright information. Also, line numbering will be removed. Journal info and pagination for accepted papers will also be assigned by the Editorial Office.

%---------
% pdftex
%---------
% The option pdftex is for use with pdfLaTeX. If eps figures are used, remove the option pdftex and use LaTeX and dvi2pdf.

%=================================================================
% MDPI internal commands
\firstpage{1} 
\makeatletter 
\setcounter{page}{\@firstpage} 
\makeatother
\pubvolume{1}
\issuenum{1}
\articlenumber{0}
%\doinum{}
\pubyear{2021}
\copyrightyear{2020}
%\externaleditor{Academic Editor: Firstname Lastname} % For journal Automation, please change Academic Editor to "Communicated by"
\datereceived{} 
\dateaccepted{} 
\datepublished{} 
\hreflink{https://doi.org/} % If needed use \linebreak

%\bibliographystyle{numeric}
\usepackage{lipsum} 

\bibstyle{numeric}
\usepackage[super]{nth}

%------------------------------------------------------------------
% The following line should be uncommented if the LaTeX file is uploaded to arXiv.org
%\pdfoutput=1

%=================================================================
% Add packages and commands here. The following packages are loaded in our class file: fontenc, inputenc, calc, indentfirst, fancyhdr, graphicx, epstopdf, lastpage, ifthen, lineno, float, amsmath, setspace, enumitem, mathpazo, booktabs, titlesec, etoolbox, tabto, xcolor, soul, multirow, microtype, tikz, totcount, changepage, paracol, attrib, upgreek, cleveref, amsthm, hyphenat, natbib, hyperref, footmisc, url, geometry, newfloat, caption

%=================================================================
%% Please use the following mathematics environments: Theorem, Lemma, Corollary, Proposition, Characterization, Property, Problem, Example, ExamplesandDefinitions, Hypothesis, Remark, Definition, Notation, Assumption
%% For proofs, please use the proof environment (the amsthm package is loaded by the MDPI class).

%=================================================================
% Full title of the paper (Capitalized)
\Title{Title}

% MDPI internal command: Title for citation in the left column
\TitleCitation{Title}

% Author Orchid ID: enter ID or remove command
\newcommand{\orcidauthorA}{0000-0000-0000-000X} % Add \orcidA{} behind the author's name
%\newcommand{\orcidauthorB}{0000-0000-0000-000X} % Add \orcidB{} behind the author's name

% Authors, for the paper (add full first names)
%\Author{Poina Lemenkova $^{1,\dagger,\ddagger}$\orcidA{}, Olivier Debeir $^{1,\ddagger}$, Thomas Lecocq $^{2,}$* and Rapha\"el De Plaen $^{1}$}
\Author{Poina Lemenkova $^{1,2}$*, Olivier Debeir $^{2}$, Thomas Lecocq $^{3}$ and Rapha\"el S. M. De Plaen $^{3}$}

% MDPI internal command: Authors, for metadata in PDF
\AuthorNames{Poina Lemenkova, Olivier Debeir, Thomas Lecocq and Rapha\"el De Plaen}

% MDPI internal command: Authors, for citation in the left column
\AuthorCitation{Lemenkova, P.; Debeir, D.; Lecocq, T., De Plaen, R.}
% If this is a Chicago style journal: Lastname, Firstname, Firstname Lastname, and Firstname Lastname.

% Affiliations / Addresses (Add [1] after \address if there is only one affiliation.)
\address{%
$^{1}$ \quad Université Libre de Bruxelles, École polytechnique de Bruxelles (Brussels Faculty of Engineering), Laboratory of Image Synthesis and Analysis (LISA). Building L, Campus de Solbosch, Avenue Franklin Roosevelt 50, Brussels 1000, Belgium. \\
$^{2}$ \quad Schmidt Institute of Physics of the Earth, Russian Academy of Sciences. Department of Natural Disasters, Anthropogenic Hazards and Seismicity of the Earth. Laboratory of Regional Geophysics and Natural Disasters (303). Moscow, 123995, Russia \\
$^{3}$ \quad Royal Observatory of Belgium, Seismology \& Gravimetry (OD2); Email: Thomas.Lecocq@seismology.be
	}

% Contact information of the corresponding author
\corres{Correspondence: polina.lemenkova@ulb.be; Tel.: +32 471 86 04 59 (P.L.)}

% Current address and/or shared authorship
% \firstnote{Current address: Affiliation 3} 
%\secondnote{These authors contributed equally to this work.}
% The commands \thirdnote{} till \eighthnote{} are available for further notes

%\simplesumm{} % Simple summary

% Abstract (Do not insert blank lines, i.e. \\) 
\abstract{\textcolor{red}{200 words max \lipsum[1]}}

% Keywords
\keyword{seismology; Galitzin seismograph; horizontal component; analogue seismogram; digitizing; earthquake recording; ground motions; historical seismograms; seismic waves} 

\begin{document}
%%%%%%%%%%%%%%%%%%%%%%%%%%%%%%%%%%%%%%%%%%

\section{Introduction}

The seismicity of the Earth represent a physical phenomena resulting from inner tectonic processes connected to energy accumulation and release \cite{GutenbergRichter}. The fundamental concepts of the Earth's physics are reflected in the movements on the surface that have different intensity of the vibrations, caused by the fluctuations of energy \cite{Shearer2019}. In the geodynamic sense, seismicity is the result of self-organisation of the Earth, responding to the internal processes of the lithosphere movements by the released energy \cite{UdiasBuforn2017}. Measuring seismic signals of the Earth by seismometers enables to evaluate ground noises and shaking of the Earth’s interior, associated with earthquakes of various magnitude. 

The importance of the seismogram datasets can be illustrated by their practical application in relation to tectonics, geology and physics of the Earth. Specifically, the wide use of seismograms includes such fundamental domains as geophysics \cite{Kulhanek}, subsurface exploration, global and planetary tectonics \cite{Crampin}, \cite{Jackson2002}, risk assessment in civil engineering \cite{Pileckietal2017}, \cite{Kurzeja}, prediction of earthquakes and relative analysis (dynamic rupture, seismic tomography and wave forms) \cite{Amorese}, volcanology \cite{Pezzo}. Analysis of seismograms is based on the interpreting of observation datasets to quantify seismic sources: location, size, rupture history, travel time and wave polarization \cite{Ammon}. Seismic data processing involves a variety of approximation methods, e.g. Fourier analysis, or an artificial modelling of the synthetic generated data \cite{Rosenbaum1994}.

For instance, earthquake-related hazards can be simulated by artificial Fortran-based stochastic time series of faults \cite{BeresnevAtkinson}. The evaluation of this method was further extended by \cite{NicknamEslamian} who generated artificial seismograms for modelling seismotectonic processes. These include, for instance, such aspects as fault length and width, analysis of sensitivity with varied epicenter locations and rupture processes showing the extent of slip of the Earth's crust, e.g. rupture velocity and shear wave velocity. Synthetic seismograms can furthermore point at the geological processes by analysis of seismic reflectors of stratigraphic data \cite{BrewMayer1998}. Besides geological domains, seismograms can be used in environmental studies to reveal the characteristics of seismic storms and their correlation with climate change and dynamic processes of the Earth \cite{Singh2021}, \cite{Toyokunietal2015}.

Predictive modelling of seismograms for analysis of earthquake hazards is largely based on the observation seismology and retrieval of information from the digital seismograms \cite{Claerbout1992}. Uniformed, standardised, converted to digital format and accurately processed seismograms enable to explain the physical nature and repeatability of earthquakes based on the ground motion signals. However, the use of such information relies on our access to the large datasets of seismograms. Generating such data pools is only possible by accurate and standardised recording system using commonly accepted and approved instruments (seismometers and seismographers) that continues through years and centuries and regularly provides data for further processing \cite{UdiasBuforn2017}. 

The history of seismography undergone a long process of technical evolution \cite{DeweyByerly1979}. In general, the instruments used for seismic data recording reflect a constant development of technologies and a state-of-the-art in methods, which provide tools to measure ground motion, earthquake's magnitude, location and frequency \cite{DeweyByerly1979}. The systematic and precise operative observation of seismicity of the Earth began in the latter part of the \nth{19} century along with a progress in technology, development of new seismographs, and general advances in geophysical methods \cite{SchweitzerLee2003}. At this time, since 1880s seismologists started recording continuously and systematically the ground motion using seismographs \cite{SchweitzerLee2003}. From that period and up to 1960s the seismograms were recorded as analogue signals using technical abilities of the narrow-band, low-range instruments of those time \cite{Batlloetal2008}.

During the first half of the \nth{20} century seismograms continued to be recorded in paper analog format \cite{Wangetal2014} or as magnetic tapes \cite{Anetal2015}. Starting with 1960s, the progress in seismography gradually continued along with the onset of computer era, development of microprocessors and advanced algorithms of geophysical data processing. Apart from advances in instrumental development, seismic data gradually increased in quantity systematically collected by seismic stations worldwide \cite{Engdahl2002}. An unprecedented initiative of global seismological network WWSSN was created in 1960s and generated a collection of high quality big seismic data \cite{OliverMurphy}, \cite{Ammon2010}. The uniformed seismic datasets of the WWSSN network were possible due to the principle of uniformity: each seismic station had identical technical equipment and a standard workflow for data calibration.

Currently, seismograms are recorded and stored digitally in high-performance data centres using advanced algorithms of computation seismology \cite{Heiner}. The progress in computational seismology now enables numerical approximations of the seismic waves for modelling signals as complex 3D structures, using advanced mathematical solutions, such as Fourier series, Chebyshev pseudospectral methods, approximation of traces by Lagrange and Chebyshev polynomials \cite{Heiner}.

However, the remaining historical datasets from previous periods form a valuable database containing archived data on seismicity of the Earth that should be used for practical scientific purposes \cite{Ambraseys1983}. The importance of these old datasets for geophysics has been discussed in relevant works \cite{Vogtetal1985}, \cite{TevesCosta1999}, \cite{Bono2007}, \cite{Okal}. Old seismograms present a key source of information regarding the seismicity of the Earth's interior up to the Middle Ages \cite{Alexandre1990a}, \cite{Alexandre1985}, \cite{Alexandre1990b}, which can be used for geophysical reconstructions  and past climate modelling. Besides, retrospective analysis of the old seismograms enables to model earthquake cycle as a continuous process and to get a better insights into the physics of the earthquakes. Finally, interpretation of the old seismograms provides information for the long-term prognosis. Predictive modelling of earthquakes is possible based on the statistical analysis of historical data series on earthquakes's locations and magnitude \cite{Kanamori1988}.

The first attempts to digitise the analogue recordings of seismic events into numeric format started in 1960s. For instance, an \emph{ad hoc} algorithm of thinning the lines on a scanned seismogram and converting the image into a 0-1 digital matrix was developed by \cite{Jarosch1975}. Another study used a device initially used for weather maps and adjusted it for automatic converting seismic signals into a digital form, and plotting them using a cathode-ray tube display  \cite{AdamsAllen1961}. The progress of digitising methods of signals accelerated in 1980s (\cite{AspinallLatchman}, \cite{CrouseMatuschka1983}) and 1990s (\cite{NorlundGriffiths1993}, \cite{McGee1995}, \cite{Harjesetal1993}) along with rapid development of computer-based algorithms and advanced geophysical software. Digitising approaches has become more and more automated since 2000s, along with a rapid development of programming languages and advanced IT tools \cite{Bietal2016}, \cite{Ishiietal2015}, \cite{Schenke}, \cite{Wangetal2014}. 
 
In recent years, the use of automated digitisers of seismograms or accelerograms (recording the acceleration of the strong motion ground) instead of manual vectorisers has received attention for the effective data processing. These include automated recognition of noise \cite{CrouseMatuschka1983}, handling smooth and wigged traces using algorithms for tracing waves based on the pixels context \cite{Bietal2016}, synchronisation of time scale for the three motion components, handling scratches, distortions or line crossings in the image, adjusting rotated position \cite{Trifunacetal1999}, to mention a few. Other problems are caused by the recording methods of the old mechanical seismographs. For instance, these include a high velocity of stylus that does not affect the paper, clipping of amplitudes, low-contrast dark background on a smoked paper, damaged records due to the bad storage conditions \cite{Baskoutasetal2000}. Digitising seismograms in such a bad conditions can be a challenge or, at least, not a straightforward task.

As a response to the arising needs for automatic, quick and accurate digitising and modelling of seismic data, various programs were developed due to the rapid growth of the programming and IT industry. Modern software for digitising seismograms are based on the principle of pattern recognition. For example, the Teseo software, developed by \cite{Pintoreetal2005}, traces raster files based on the piecewise cubic Bézier curves. This vectoriser handles the monochrome images by three methods: manual, automatic, and neural networks. A software program using C\# is presented by \cite{Wangetal2016} using algorithm of color scene recognition by Color Scene Filed Method. An interactive program using Matlab GUI was developed by \cite{XuXu} using an algorithm of inversion problem for a matrix of model parameters and the observed seismic data. 

Another example of seismic software, the The SeisDig, was developed by Berkeley Seismological Laboratory under a copyright. It presents an interactive Matlab GUI based tool aimed at digitising scanned images, performing trace inspection, checking time conservation, recognising individual events, and inputing metadata as a header \cite{BromirskiChuang2003}. An innovative approach to the development of the seismic recorder is presented by \cite{Khanetal2012} who developed the DigiSeis software. The DigiSeis is aimed at digitizing seismic signals using a built-in PC sound card for processing seismic signals. It visualises the two channel seismic data in time and frequency, and enables to perform a time series analysis. A special attempt for automatic extracting of wave data from paper seismograms was presented by \cite{Liuetal2001} with a developed seismogram digitization and DB management system using Delphi3 based on the syntax of Pascal language. 
 
%%%%%%%%%%%%%%%%%%%%%%%%%%%%%%%%%%%%%%%%%%
\section{Materials and Methods}

\subsection{Data}
The data were obtained from the archives of Royal Observatory of Belgium, Department of Seismology \& Gravimetry, the Uccle seismic station of Belgium (UCC). Established in 1899 by Eugène Lagrange, Uccle was the only seismic station in Belgium for the most of \nth{20} century. Afterwards seismic monitoring  gradually expanded and resulted in large seismic network system of Belgium with Uccle still being its important centre. The dataset received from the Uccle station presents a large collection of scanned paper-based seismograms showing the set of registered waveforms of seismic swarm that occurred between 1953 and 1957 close to the town of Court-Saint-Etienne, 20 km SE of Brussels. Similar seismic storm occurred recently (2008-2010) and was interpreted by \cite{VanNotenetal2015} as an evidence of fault processes in a weakened crust within the Cambrian core of the Brabant Massif, Belgium.

The seismograms have been receive using the Galitzine seismometer. This type of seismographs is constructed using a following principle: a coil of wire is fixed to the pendulum which oscillates between the poles driven by the earthquake forces. The movement of coil between a strong magnet generate an electric induction current. The current is then being transferred to a sensitive galvanometer with signals registered photographically on the seismogram \cite{Neumann}. It is the \nth{1} electromagnetic instrument developed by B.B. Galitzine in 1910, designed to record a single component (horizontal or vertical) for earthquakes recording \cite{Wgg}. 

B.B. Galitzine reduced the sensitivity of the system by diverting a part of the current and thus improved the accuracy of the system. The special requirements for standard type of operation of these units include the following: (1) Damping the resistance to the galvanometer (2) Damping the seismometer pendulum through the cumulated effect of a magnetic damping device and the external resistance to the pendulum coil. These conditions are satisfied by the Galitzine system that includes a copper plate, fixed on the rod next to the coil, which oscillates in the field of a second magnet and allows for damping by Foucault current. Due to such effectiveness, the Galitzine seismometer is still being used for seismic measurements of Galitzine type \cite{ChakrabartyTandon}.

In this study, for processing seismograms we used the DigitSeis software for processing analog scanned seismograms \cite{BogiatzisIshii2016}.


%%%%%%%%%%%%%%%%%%%%%%%%%%%%%%%%%%%%%%%%%%
\section{Results}

\subsection{Subsection}
\subsubsection{Subsubsection}

\textcolor{red}{\lipsum[1]}

\subsection{Subsection}

\textcolor{red}{\lipsum[1]}

\subsection{Subsection}

\textcolor{red}{\lipsum[1]}

%%%%%%%%%%%%%%%%%%%%%%%%%%%%%%%%%%%%%%%%%%
\section{Discussion}

\textcolor{red}{\lipsum[1]}

%%%%%%%%%%%%%%%%%%%%%%%%%%%%%%%%%%%%%%%%%%
\section{Conclusions}

\textcolor{red}{\lipsum[1]}

%%%%%%%%%%%%%%%%%%%%%%%%%%%%%%%%%%%%%%%%%%
\vspace{6pt} 

%%%%%%%%%%%%%%%%%%%%%%%%%%%%%%%%%%%%%%%%%%
\authorcontributions{Conceptualization, O.D. and P.L.; methodology, O.D. and P.L.; software, P.L.; validation, O.D., S.L., R.D.P. and P.L.; formal analysis, O.D., S.L., R.D.P. and P.L.; investigation, O.D., S.L., R.D.P. and P.L.; resources, R.D.P..; data curation, R.D.P.; writing---original draft preparation, P.L.; writing---review and editing, O.D., S.L. and R.D.P.; visualization, P.L.; supervision, O.D. and T.L.; project administration, O.D.; funding acquisition, O.D., according to the  \href{http://img.mdpi.org/data/contributor-role-instruction.pdf}{CRediT taxonomy} of terms. All authors have read and agreed to the published version of the manuscript.}

\funding{This research received funding from the Université Libre de Bruxelles, École polytechnique de Bruxelles (Brussels Faculty of Engineering), Laboratory of Image Synthesis and Analysis (LISA), Brussels, Belgium.}

\dataavailability{Data supporting reported results are the courtesy of Royal Observatory of Belgium, Department of Seismology \& Gravimetry (OD2), used in this study as available archived datasets of seismograms.}

\acknowledgments{The authors thank the anonymous reviewers for their comments and suggestions on the initial version of the manuscript.}

\conflictsofinterest{The authors declare no conflict of interest neither any personal circumstances or interest that may be perceived as inappropriately influencing the representation or interpretation of reported research results.} 

%%%%%%%%%%%%%%%%%%%%%%%%%%%%%%%%%%%%%%%%%%
%% Optional
\abbreviations{Abbreviations}{
The following abbreviations are used in this manuscript:\\

\noindent 
\begin{tabular}{@{}ll}
	WWSSN & World-Wide Standardized Seismograph Network \\
\end{tabular}}

%%%%%%%%%%%%%%%%%%%%%%%%%%%%%%%%%%%%%%%%%%
\end{paracol}

\reftitle{References}

%=====================================
% References, variant A: external bibliography
%=====================================

\externalbibliography{yes}
\bibliography{Seis.bib}
%\nocite{*}

% If authors have biography, please use the format below
%\section*{Short Biography of Authors}
%\bio
%{\raisebox{-0.35cm}{\includegraphics[width=3.5cm,height=5.3cm,clip,keepaspectratio]{Definitions/author1.pdf}}}
%{\textbf{Firstname Lastname} Biography of first author}
%
%\bio
%{\raisebox{-0.35cm}{\includegraphics[width=3.5cm,height=5.3cm,clip,keepaspectratio]{Definitions/author2.jpg}}}
%{\textbf{Firstname Lastname} Biography of second author}

% The following MDPI journals use author-date citation: Admsci,  Arts, Econometrics, Economies, Genealogy, Humanities, IJFS, Jintelligence, JRFM, Languages, Laws, Literature, Religions, Risks, Social Sciences. For those journals, please follow the formatting guidelines on http://www.mdpi.com/authors/references
% To cite two works by the same author: \citeauthor{ref-journal-1a} (\citeyear{ref-journal-1a}, \citeyear{ref-journal-1b}). This produces: Whittaker (1967, 1975)

%%%%%%%%%%%%%%%%%%%%%%%%%%%%%%%%%%%%%%%%%%
\end{document}

